\chapter{A short proof for every noncommuting residue}
\label{component:p2c3fc22497a5}
\noindent\textbf{Component }\path{p2c3fc22497a5}\par
{\small\noindent Source: \path{CORE/prior/inherited/penpaper_noncommuting.tex}\par}
{\small\noindent SHA256: \path{2c3fc22497a59086acf6bc93ff5f5408d84300149843aa9d4f6239fb6847d29a}\par}
{\small\noindent Inventory use: full exposition | entire noncommuting-residue proof; full exposition | C14\_\allowbreak{}ENTIRE\_\allowbreak{}ORIGINAL\_\allowbreak{}PARENT\_\allowbreak{}GRAPH; precise subordinate uses retained below\par}
\medskip
\section{A short proof for every noncommuting residue}
\label{p2c3fc22497a5:sec:penpaper-noncommuting}

The noncommuting part of the argument has a simple organizing principle:
\emph{a vanishing event constrains an entangled state, and there are only four
compatible such constraints}.  Coefficients can occur at arbitrarily many
relative scales, but the proof only changes when a new independent constraint
appears.  The number of relevant changes is therefore bounded by geometry.

This section gives a short reading route through
Theorem~\ref{p94dfbc8a5970:thm:all-noncommuting}.  Its conclusion concerns independent tests
whose two measurement pairs are noncommuting \emph{in residue}.  It does not
include infinitesimal approaches to commuting measurements, and it is not by
itself the global finite SDP theorem.

\subsection{What must be certified}
Keep the coefficient field $K_c$ and the testing field $K_t$ separate, and put
$T=K_c\otimes_{\mathbb R}K_t$.  A source is physical when its Bell slack is
nonnegative at every experiment over $K_c$.  The assertion to prove is stronger
than that ordinary positivity: its evaluation at an independent experiment
over $K_t$ must lie in $\Sigma T^2$.

Write the entire parent in the form
\begin{equation}\label{p2c3fc22497a5:eq:pp-packets}
 F=\sum_{e}c_e\,b_e(u,v)R_e(W,X,Y,Z;u,v)^2+d\,C,
 \qquad c_e\geq0,\quad d\geq0.
\end{equation}
Here $C=D_uD_v\langle H\rangle$, with the state norm left homogeneous;
the factors $b_e$ are the positive testing factors in
\eqref{p0c8f420148b7:eq:parent}, including its factor four.  The $R_e$ are the eight
measurement amplitudes.  For example,
\[
 R_n=X,\quad R_A=Y,\quad
 R_m=vW-X+uvY-uZ,\quad
 R_D=uW+uvX-Y-vZ.
\]
Normalize $d+\sum_e c_e=1$, unless the source is zero.  Choose a bounded
projective state chart and finite half angles whose residues satisfy
$u_0v_0\ne0$.  These choices use only pure testing factors.

Only one algebraic certificate is needed throughout.  If a finite identity
has the form
\begin{equation}\label{p2c3fc22497a5:eq:pp-strict}
 F/\sigma=a+h^TQh,\qquad a\in K_c,\quad a\geq0,
 \quad Q\in\operatorname{Sym}(B_c\otimes B_t),\quad\bar Q\succ0,
\end{equation}
where $\sigma>0$ is pure coefficient, then the strict Gram lemma gives
$F\in\Sigma T^2$.  The entries of $h$ need not be bounded.  In particular,
rescaling a normal coordinate by a large \emph{pure coefficient} is allowed.
The argument never takes an inverse or a square root of a mixed tensor.

There are two immediate cases.  Positive slack residue is covered by
Lemma~\ref{p25af2c58da7e:lem:residue}.  Literal $d=0$ makes
\eqref{p2c3fc22497a5:eq:pp-packets} an event-square sum.  All subsequent steps keep every
packet of the complete parent, including its infinitesimal terms.

\subsection{Entanglement turns the events into a matching}
For a product state, the CHSH expectation obeys
\[
 S=a_0(b_0-b_1)+a_1(b_0+b_1)
 \leq |b_0-b_1|+|b_0+b_1|=2\max(|b_0|,|b_1|).
\]
At noncommuting residue angles equality $S=2$ is impossible.  If both Bob
moduli are below one it is already strict.  If one equals one, the other is
strictly smaller, and equality would require a common eigenstate of Alice's
two noncommuting reflections.  Thus the signed core has positive testing
residue at every product residue state.  Its actual testing value is positive,
so $dC$ is a square after separate field square roots.  Such tests are settled.

Now suppose $\bar d=0$, the slack residue vanishes, and the residue state is
entangled.  Write that real state as an invertible $2\times2$ matrix $\Psi$.
For any fixed Alice ray $a$, the equation
\[
                         a^T\Psi b=0
\]
selects exactly one Bob projective ray, and conversely.  Two vanishing events
cannot therefore share a local outcome ray.  The losing events occur in a single bipartite cycle of local outcome rays,
with cyclic edge order
\[
                    (\lambda,D,c,A,\alpha,m,\beta,n).
\]
Consequently every set of vanishing events is a matching, of size at most
four.  This is the finiteness mechanism.

Any three matching event vectors are state-linearly independent.  Indeed,
the span of two rank-one matrices with independent left and right factors
contains no third rank-one projective point.  A dependence among three
rank-one matrices must therefore have a common left or right factor;
matching edges cannot do so.  For at most three vanishing events, their
amplitudes may hence be used as exact local state coordinates.  A nonzero
minor at the chosen residue supplies the inverse change of coordinates;
its denominator is a pure testing unit.  This is an exact coordinate change,
not a linearization.

Four matching edges give one of two perfect matchings.  In the representative
$\{n,m,A,D\}$, the equations are
\[
 X=Y=0,\qquad vW=uZ,\qquad uW=vZ.
\]
Thus $u^2=v^2$.  Choosing its sign and putting $W=1$ gives the smooth curve
$v=u$, $Z=1$.  It has four independent normal directions: three state
directions and the angle difference.  The other perfect matching is its
CHSH-symmetric counterpart.  Section~\ref{p76173e99815f:sec:zero-d-cycle} establishes these
statements directly from the event amplitudes.

\subsection{Unequal scales are handled one normal coordinate at a time}
Suppose $k\leq3$ event weights $a_1,\ldots,a_k$ have already been identified,
and every remaining coefficient is infinitesimal relative to each $a_i$.
Let $\sigma$ be the largest remaining coefficient, including $d$, and put
\[
 \varepsilon_i=\sigma/a_i\in\mathfrak m_c.
\]
In exact event coordinates the parent has the form
\begin{equation}\label{p2c3fc22497a5:eq:pp-diagonal}
 F(x,y)=\sum_{i=1}^k a_i b_i(y)x_i^2+\sigma G(x,y),
 \qquad \bar b_i(y_0)>0,
\end{equation}
with bounded coefficients after clearing pure testing unit denominators.
The variables $y$ comprise the remaining projective state and angle
coordinates.  On the exact hard kernel $x=0$, physicality implies
\[
                         g(y):=\overline G(0,y)\geq0.
\]

If $g(y_0)>0$, the roots $(x_i/\sqrt{\varepsilon_i},1)$ immediately give a
strict residue Gram.  If $g(y_0)=0$ and $\nabla^2g(y_0)\succ0$, solve for an
exact coefficient-field stationary center.  The correct scaling is
$x_i=\varepsilon_i s_i$, separately for every normal.  Dividing the stationary
equations gives
\begin{align*}
 2b_i(y)s_i+G_{x_i}(\varepsilon s,y)&=0,\\
 \sum_i\varepsilon_i s_i^2\nabla b_i(y)
                 +\nabla_yG(\varepsilon s,y)&=0.
\end{align*}
Their residue Jacobian is block upper triangular, with diagonal blocks
\[
 2\operatorname{diag}(\bar b_i(y_0)),\qquad \nabla^2g(y_0).
\]
Both are positive definite.  The semialgebraic implicit-function theorem,
transferred to $K_c$, gives $(x_c,y_c)$ with
$x_{ci}=\varepsilon_i s_{ci}$ and bounded $s_c$.

Finite Taylor expansion at this exact center, on the roots
\[
 h=\big((x_i-x_{ci})/\sqrt{\varepsilon_i}\big)_{i=1}^k
                  \mathbin{\oplus}(y-y_c),
\]
has residue Gram
\[
 \bar Q=\operatorname{diag}
       \big(\bar b_1,\ldots,\bar b_k,\tfrac12\nabla^2g(y_0)\big)\succ0.
\]
Hard normal--tangent terms acquire a factor $\sqrt{\varepsilon_i}$;
hard pure tangent terms acquire $\varepsilon_i$.  Higher Taylor terms leave
an actual infinitesimal deviation in their coefficient.  The constant
$F(x_c,y_c)/\sigma$ is nonnegative by physicality, giving
\eqref{p2c3fc22497a5:eq:pp-strict}.  This proves Lemma~\ref{p94dfbc8a5970:lem:diagonal-hard-scales} and
explains why no comparability of the $a_i$ is needed.

The implicit-function step is local in a fixed finite-dimensional real
parameter space.  Its existence and uniqueness statements transfer to the
coefficient field.  It is not an infinite asymptotic expansion or an
assumption of compactness over a non-Archimedean field.

\subsection{When the signed term appears, physicality forces curvature}
Suppose $\overline{d/\sigma}>0$ at the next scale.  With no hard constraint,
or with one or two hard event constraints, the constrained state form is
strictly positive on its subspace $Z=0$.  At a zero it therefore has a simple
kernel and positive complementary state pivots.

The remaining angular curvature is encoded by one small determinant.
The parent determinant is a positive angular factor times $UVL(U,V)$,
where $U=u^2$, $V=v^2$, and
\begin{equation}\label{p2c3fc22497a5:eq:pp-laurent}
 L(U,V)=h+\sum_{e\in E}q_eU^{e_1}V^{e_2},
 \qquad q_e>0,\qquad
 E\subset\{-1,0,1\}^2\setminus\{(0,0)\}.
\end{equation}
All outer coefficients are nonnegative, as follows from the factored
determinant formulas in Section~\ref{p0c8f420148b7:sec:physical}.  Constraining one or two
event amplitudes amounts to taking the coefficient of their weights in
this determinant: in an adapted state basis it is the determinant of the
remaining state form, times a positive factor.  Hence it retains the same
Laurent structure.

At a physical zero, logarithmic stationarity and curvature are exactly
\begin{equation}\label{p2c3fc22497a5:eq:pp-newton-hessian}
 \sum_{e\in E}w_e e=0,\qquad
 \nabla^2_{\log U,\log V}L=\sum_{e\in E}w_e ee^T,
 \qquad w_e=q_eU^{e_1}V^{e_2}>0.
\end{equation}
Thus strict angular curvature is equivalent to the active exponent vectors
spanning the plane.  There is no eigenvalue estimate to compute.

For clarity, the finite support argument establishing that span has only
two possible obstructions.  A stationary support of rank one must lie on
one of the four opposite-pair lines of the Newton square.  Substituting the
factored coefficients shows that it either activates an exponent off that
line, or leaves a perfect matching of events.  Rank zero leaves a support
contained in a nonextendable matching triple.  For the unconstrained rank-zero case, the forbidden pairs form this
event cycle together with its four opposite chords.  An independent triple
must have cyclic gaps $(2,3,3)$: the only other possible gap pattern
$(2,2,4)$ contains an opposite pair.  Rotation of the cycle carries all such
triples to $\{\lambda,m,c\}$.  Indeed the CHSH maps $g_T$ and $g_B$ act on
edge indices as $i\mapsto2-i$ and $i\mapsto3-i$ modulo eight, so
$g_Bg_T$ is a one-edge rotation.  The representative has the exact kernel state
\[
 \psi_H=(uv,-u,u^2v,1+v^2+u^2v^2),\qquad
 C(\psi_H)=-2D_uD_vu^2v^2<0.
\]
The rank-one perfect-matching case has a
matching curve on which the core is zero but its differential is nonzero
in every event-normal direction.  In the representative matching, with
$B=1+t^2$, that differential is
\[
              2t\big(B\,dX-B\,dY-dR_m+dR_D\big).
\]
Imposing at most two event constraints leaves a direction with nonzero
core derivative.  All event squares have zero derivative there, so the
constrained parent would change sign.  Either obstruction contradicts
physicality.  The short exact support arguments are proved in
Theorem~\ref{pe25c50379365:thm:curvature} and
Propositions~\ref{pb16e5066334b:prop:one-hard-event}--\ref{pf47f6a4e3a60:prop:two-hard-event}; they require
only the displayed determinant factors and these kernel witnesses.

For three hard events there is no state variable left, and there are only
two matching orbits.  Their exact kernel states can be chosen as
\[
 \psi_N=(uv,-u,0,1+v^2),\qquad \psi_E=(u,0,0,v).
\]
After division by positive testing factors, their restrictions are
\begin{align*}
 L_N={}&\frac{2n(1+U)}V+2cU+\frac{2D(1+V)}U
       +\frac{2\alpha(1+U)(1+V)}{UV}+2\beta V-d,\\
 L_E={}&d(U/V-1)+\frac{2\lambda(1+U)}V
       +\frac{2cU(1+V)}V+\frac{2D(U-V)^2}{UV}
       +\frac{2\alpha(1+U)}U+2\beta(1+V).
\end{align*}
For $L_N$, a positive $n,D$ or $\alpha$ already supplies two independent
exponents; if they all vanish, the remaining positive exponentials cannot
have a stationary zero.  For $L_E$, $d>0$ supplies $(1,-1)$.  Failure of
two-dimensional span forces $\lambda=c=\alpha=\beta=0$; the remaining
one-variable expression has minimum
\[
 2\sqrt{2D(d+2D)}-d-4D<0,
 \qquad (d+4D)^2-8D(d+2D)=d^2>0.
\]
Again physicality excludes the obstruction.  Multiplication by a positive
factor preserves the Hessian sign at a stationary zero.  The positive
state pivots, when present, then promote angular curvature to the full
constrained Hessian by an ordinary real Schur complement.  This Schur
complement is used only in the residue field.  The weighted argument above
now supplies the exact tensor certificate.

\subsection{Four constraints leave one Laurent variable}
It remains to understand what happens if all four matching events appear
before the signed term does.  Take the matching $\{n,m,A,D\}$ and, by
symmetry, choose its smallest weight to be $D=\mu$.  The hypotheses are
\[
 d/\mu\in\mathfrak m_c,\qquad
 \lambda/\mu,c/\mu,\alpha/\mu,\beta/\mu\in\mathfrak m_c.
\]
Set $\eta_1=\mu/n$, $\eta_2=\mu/A$, $\eta_3=\mu/m$, so
$0<\eta_i\leq1$; their residues may be zero.  Near the matching curve use
$W=1$ and exact coordinates
\[
 R=v-X+uvY-uZ,\qquad
 Z=(v-X+uvY-R)/u,\qquad \delta=v-u.
\]
The fourth amplitude equals $f+L$, where
\[
 f=(u^2-v^2)/u,\qquad
 L=v(u+u^{-1})X-(1+v^2)Y+(v/u)R.
\]
On $X=Y=R=0$ the signed-plus-fourth packets are exactly
\begin{equation}\label{p2c3fc22497a5:eq:pp-four-completion}
 4\mu f^2+2duf
 =4\mu\left(f+\frac{du}{4\mu}\right)^2
                  -\frac{d^2u^2}{4\mu}.
\end{equation}
Let $\tau=\max(\lambda,c,\alpha,\beta)$.  Testing the negative minimum in
\eqref{p2c3fc22497a5:eq:pp-four-completion} at $u=1$ proves
\[
                  \tau>0,\qquad \rho=\frac{d^2}{\mu\tau}
                  \text{ is bounded}.
\]
Thus physicality itself bounds the first adverse normal correction by the
complementary event scale.

Put $\varepsilon=d/\mu$ and scale the four stationary normal coordinates by
\[
 X=\varepsilon\eta_1s_1,\quad
 Y=\varepsilon\eta_2s_2,\quad
 R=\varepsilon\eta_3s_3,\quad \delta=\varepsilon r,
 \qquad u=t.
\]
For $\bar\rho>0$, divide their stationary equations by $d$.  Writing
$B=1+t^2$, $\ell=(B,-B,1)$, $b=(B^2,B^2,1)$ and
$k=-2r+\sum_i\bar\eta_i\ell_i s_i$, their residues are
\[
 8b_i s_i+8k\ell_i+(4tB,-4tB,0)_i=0,
 \qquad -16k-4t=0.
\]
Subtracting suitable multiples of the last row leaves positive diagonal
normal pivots $8b_i$ and $32$.  In particular the system remains nonsingular
when any $\bar\eta_i$ vanishes.  Its solution is
\[
 s_1=-\frac{t}{4B},\quad s_2=\frac{t}{4B},\quad s_3=\frac t4,
 \quad r=\frac t8(1-\bar\eta_1-\bar\eta_2+\bar\eta_3).
\]
After this normal completion, the residue potential along the matching
curve is exactly
\[
 g(t)=4(1+t^2)(P+Qt^2)
       -\frac{\bar\rho t^2}{4}(1+\bar\eta_1+\bar\eta_2+\bar\eta_3),
\]
where $P=\overline{(\lambda+\alpha)/\tau}$,
$Q=\overline{(c+\beta)/\tau}$ and $P+Q>0$.  Its quotient by $4t^2$ is
\begin{equation}\label{p2c3fc22497a5:eq:pp-final-circuit}
 P(1+t^{-2})+Q(1+t^2)
       -\frac{\bar\rho}{16}(1+\bar\eta_1+\bar\eta_2+\bar\eta_3).
\end{equation}
Physicality makes this nonnegative.  At a zero, stationarity forces
$P/t^2=Qt^2$, so $P,Q>0$ and the logarithmic second derivative is positive.
There is therefore a strict tangent pivot as well as the four strict normal
pivots.  If $\bar\rho=0$, the complementary event value is already strictly
positive and the direct normal Gram suffices.

For the exact restoration, the roots at a coefficient-field stationary
center are
\[
 \left(\sqrt{n/\tau}(X-X_c),\sqrt{A/\tau}(Y-Y_c),
 \sqrt{m/\tau}(R-R_c),\sqrt{\mu/\tau}(\delta-\delta_c),t-t_c\right).
\]
The key bounds are transparent: for a normal weight $a_i=\mu/\eta_i$,
\[
 \frac{a_ix_{ci}^2}{\tau}=\rho\eta_i s_{ci}^2,
 \qquad
 \frac{a_ix_{ci}}{\sqrt{a_i\tau}}=\sqrt{\rho\eta_i}\,s_{ci}.
\]
Every quantity on the right is bounded.  Exact packetwise Taylor allocation
therefore gives a strict Gram, whose tangent Schur complement is $g''/2$.
Theorem~\ref{pa2cbe6ecbeb4:thm:matching-flag-terminal} supplies the full details.  No inverse
of $\eta_i$ is needed in a bounded Gram coefficient.

\subsection{The finite induction}
Start with all event coefficients having positive residue.  If the residue
slack is zero and $\bar d=0$, they form a nonempty matching $S$ at the
entangled residue state.  For $|S|\leq3$, take the next scale
\[
             \sigma=\max\bigl(\{d\}\cup\{c_e:e\notin S\}\bigr).
\]
Every weight already in $S$ dominates $\sigma$.  Restricting to their exact
event kernel gives the next-scale nonnegative residue $g$.

If $g$ is positive, the weighted certificate finishes.  If $g=0$ and
$d/\sigma$ has positive residue, the Laurent curvature argument finishes.
Otherwise $g$ is an event-square sum.  At least one remaining coefficient
has positive residue after division by $\sigma$, and every such event must
vanish at the same entangled residue state.  Add all of them to $S$.
The enlarged set is still a matching, and every coefficient outside it is
infinitesimal relative to each of its weights.

Each repetition adds an edge.  A matching has at most four edges, so there
are at most three enlargements.  On reaching four edges the preceding
one-variable terminal applies, independently of their relative weights.
If $\bar d>0$ at the outset, the unconstrained Laurent curvature theorem
and the regular-contact certificate already apply.  These cases exhaust
all noncommuting residue tests and prove Theorem~\ref{p94dfbc8a5970:thm:all-noncommuting}.

The insight is that the proof counts constraints, not coefficient regimes.
The event cycle supplies a finite depth; physicality supplies curvature at
the first visible signed scale; exact Taylor expansion converts that
curvature into a separated-field sum of squares.  The commuting boundary
requires additional arguments because the outcome rays then coalesce and
the matching argument no longer controls the normal directions.
